\section{Complete Practical Coding}
\label{sec:practical-03}

Use semantic commands for meaning and keep manual styling restrained.

\begin{latexcode}{Complete practical file: readable scientific prose}
\documentclass{article}
\usepackage[T1]{fontenc}
\usepackage{lmodern}
\usepackage[british]{babel}

\newcommand{\species}[1]{\textit{#1}}

\begin{document}

\section{Observation}
We observed \species{Example species} in 12 plots. The first
measurement was made in the morning; the second was made two
hours later. The difference was small---not zero---and should be
reported without unnecessary emphasis.

The sample identifier was \texttt{plot\_07}. A short footnote can
hold a useful qualification without interrupting the sentence.\footnote{The
demonstration values are not research data.}

\section{Interpretation}
Use \emph{emphasis} for meaning. Use \textbf{bold text} mainly for
short labels, not whole paragraphs.

\end{document}
\end{latexcode}

\begin{codelab}
Replace both paragraphs with your own prose, add one species name, and confirm
that every identifier and reserved character prints correctly.
\end{codelab}
